\documentclass[11pt,letterpaper]{article}
\usepackage[margin=0.85in]{geometry}
\usepackage{lmodern}
\usepackage{hyperref}
\hypersetup{colorlinks=false, urlcolor=black, pdfborder=0 0 0}
\setlength{\parindent}{0pt}
\setlength{\parskip}{0.45em}
\raggedright
\pagestyle{empty}

\begin{document}

\noindent
{\bfseries Ke-Cheng Hsu (Clio)}\\
Ph.D.\ Candidate, Department of Economics, HKUST\\
(852) 6840\,9753 \quad $|$ \quad \href{mailto:khsu@connect.ust.hk}{khsu@connect.ust.hk}\\
\href{https://cliohsu1997.github.io/}{cliohsu1997.github.io} \quad $|$ \quad \href{https://github.com/cliohsu1997}{github.com/cliohsu1997}

\vspace{0.7em}
\noindent
October 21, 2026

\vspace{0.35em}
\noindent
Hiring Committee\\
Compass Lexecon --- Beijing / Shanghai\\
Re: Economist --- Beijing --- Q3 2026 (JR260969)

\vspace{0.2em}
Dear Hiring Committee,

I am writing to apply for the Economist position in Compass Lexecon's Beijing/Shanghai office. I am a Ph.D.\ candidate in Economics at HKUST specializing in empirical industrial organization and digital economics. I am native in Mandarin and professionally fluent in English, and I am eager to apply structural and reduced-form tools to competition, regulation, and damages matters in Asia-Pacific.

My research centers on using large proprietary datasets to analyze welfare and conduct counterfactual analysis. In ``Too Much Unsold Stock?: Environmental Consequences of Price Commitment,'' joint with Kei Kawai (Senior Consultant, Compass Lexecon) and Kohei Kawaguchi (Advisor, UTokyo Economic Consulting Inc.), we study whether resale price maintenance can be pro-competitive by encouraging sellers to hold inventory, and how much of the resulting unsold stock is an environmental cost of price commitment. Estimating a dynamic oligopoly model on monthly bookstore data, we find that marginal-cost estimates differ meaningfully: the classical static model overestimates marginal cost by \(\sim\)20\% when ignoring dynamic demand uncertainty and inventory decisions. Through this collaboration, I learned to communicate clearly and rigorously with data: cleaning and transforming raw records, then using descriptive statistics, regression, and data visualization to convey the central idea before turning to structural models, and to explain key intuition with simple, easy-to-understand graphs.

My job market paper studies return policy as a trade-off between purchase and return efficiency: a stricter return policy curbs inefficient ex-post returns, but the resulting lower return rates make higher prices more attractive for sellers and can exclude ex-ante efficient purchases. Empirically, automatic provision of return-shipping insurance performs better in high-margin categories by expanding valuable purchases, while selective adoption of return-shipping insurance outperforms automatic provision by \(\sim\)15--30 percentage points of first-best welfare in thinner-margin categories by curbing excessive returns; a partial return-penalty discount (discounted price + \(\sim\)10--15\% restocking fee) beats both outside dresses, attaining \(\sim\)75--78\% of first-best welfare. I build reproducible pipelines on AWS and HPC/SLURM in Python, R, and C++, and write tests to ensure clarity and rigor; I am comfortable with GMM, instrumental variables, bootstrap inference, and counterfactual welfare analysis. In Ph.D.-level industrial organization coursework, I also conducted merger simulation, market entry/exit counterfactuals, and identification of firm conduct parameters.

Although my CV lists an expected graduation of June 2027, I can accelerate completion and join as early as Q4 2026 if an offer is available. I have attached my CV and transcript, and I can provide a writing sample (job market paper) through the portal. Thank you for your consideration.

\vspace{0.7em}
Sincerely,\\[1.0em]
Ke-Cheng Hsu (Clio)

\end{document}
